\documentclass[10pt,a4paper]{amsart}
\usepackage[margin=19mm]{geometry}
\usepackage{amsmath,amssymb,mathtools}
\usepackage[scr=rsfs,cal=euler]{mathalfa}
\usepackage{xfrac}
\usepackage[unicode,hidelinks,bookmarks=false]{hyperref}
\newcommand{\ie}{i.e.\ }
\newcommand{\cf}{cf.\ }
\newcommand{\resp}{respectively\ }
\newcommand{\Subsection}{\subsection}
\providecommand{\nobreakdash}{\nobreak\hbox{-}}
\setlength{\emergencystretch}{2em}
\multlinegap0pt
\usepackage{bm}
\usepackage{mathtools}
\usepackage{xcolor}
\usepackage{float}
\usepackage{siunitx}
\usepackage{tcolorbox}
\newtheorem{lemma}{Lemma}
\usepackage{cleveref}
\crefname{lemma}{Lemma}{Lemmas}
\newcommand{\defeq}{\vcentcolon=}
\newcommand{\hquad}{\hspace{0.5em}}
\DeclareMathOperator*{\minimize}{minimize}
\newcommand{\norm}[1]{\left\lVert#1\right\rVert}
\DeclareMathOperator*{\st}{subject\ to}
\title{Deep FlexQP: Accelerated Nonlinear\\Programming via Deep Unfolding}
\author{Alex Oshin, Rahul Vodeb Ghosh, Augustinos D. Saravanos, Evangelos A. Theodorou}
\date{Source: arXiv:2512.01565v2 (CC BY 4.0)}
\begin{document}
\maketitle

\noindent\emph{Source and licence.} Deep FlexQP: Accelerated Nonlinear\\Programming via Deep Unfolding. Version arXiv:2512.01565v2 (CC BY 4.0), distributed by arXiv under the Creative Commons Attribution 4.0 International licence, http://creativecommons.org/licenses/by/4.0/, as recorded in the arXiv licence metadata for this exact version and shown on the abstract page https://arxiv.org/abs/2512.01565v2. Full licence notice: \texttt{licenses/arxiv-2512.01565-CC-BY-4.0.txt}.\par\smallskip

\noindent\textit{Editorial scope.} Counted unit: the article's lemma eq:relaxed\_qp\_alt (source0.tex lines 1105--1110). The article's complete original proof of it is delivered with it (source0.tex lines 1111--1113, one \texttt{proof} environment of 3 lines). No mathematics is written, repaired or re-derived: the statement and the proof are the article's own lines, in the article's own notation, and no retained line is substituted or re-flowed.  Retained uncounted as the source-local context the proof needs: lines 688--690; lines 710--712.  Nothing was omitted from any retained span, so the package carries no omission record.  No retained line contains a citation, so the wrapper carries no bibliography and no bibliography entry.  No retained line contains a figure environment, an image-inclusion command or drawing code, so no image is delivered and the package declares no asset dependency.  Editorial formatting only, in addition: an AMS \texttt{amsart} preamble that restates the article's own declarations and macros used by the retained lines, namely the environments \texttt{lemma} on separate counters, together with the standard packages those lines need.  The retained environments print in this document as Lemma 1 in order of appearance, whereas the article prints the same items with the numbering of its own full document.  The complete native source of the article as distributed in the arXiv e-print for this exact version is bundled unchanged as \texttt{sources/190131/source0.tex}.
\begin{equation}
    \minimize_x \hquad   \frac{1}{2} x^\top P x + q^\top x, \quad \st \hquad Gx \leq h, \hquad Ax = b, \label{eq:qp}
\end{equation}

\begin{equation}
    \minimize_{x, s \geq 0} \hquad \phi(x, s; \mu_I, \mu_E) \defeq \frac{1}{2} x^\top P x + q^\top x + \mu_I \norm{Gx + s - h}_1 + \mu_E \norm{Ax - b}_1, \label{eq:relaxed_qp}
\end{equation}

\begin{lemma}
The relaxed QP \cref{eq:relaxed_qp} can equivalently be expressed as the following optimization:
\begin{equation} \label{eq:relaxed_qp_alt}
    \minimize_x \quad \phi(x; \mu_I, \mu_E) \defeq \frac{1}{2} x^\top P x + q^\top x + \mu_I \norm{(Gx - h)_+}_1 + \mu_E \norm{Ax - b}_1.
\end{equation}
\end{lemma}

\begin{proof}
    Convert the optimization to the form without slack variables. This can be equivalently derived by directly relaxing the constraints of the original QP \cref{eq:qp} using $\ell_1$ penalties.
\end{proof}

\end{document}
